\section{Zeitbereichsanalyse von $R$-$L$-Lasten bei Schaltbetrieb}

Diese Section behandelt den \textbf{Kernfall der Leistungselektronik}:
Ein $R$-$L$-Lastkreis, der durch einen idealen Schalter periodisch ein- und ausgeschaltet wird.
Ziel ist die \textbf{exakte Zeitbereichslösung} des Stroms und der Spannungen in den einzelnen Schaltintervallen.

\smallskip

\textbf{Business-Relevanz:}
\begin{itemize}
    \item Dimensionierung von Induktivität und Schaltfrequenz
    \item Berechnung von Stromrippel und Spitzenstrom
    \item Stationärer periodischer Betrieb
    \item Grundlage für Wirkungsgrad, Verluste und Bauteilauslegung
\end{itemize}

\subsection{Grundschaltung und Topologien}

\subsubsection{Ideale Grundschaltung}

\begin{center}
\begin{tikzpicture}[scale=1]
\draw (0,0) to[battery,l=$U_{\mathrm{DC}}$] (0,3);
\draw (0,3) -- (2,3);
\draw (2,3) to[switch,l=$S$] (4,3);
\draw (4,3) to[R,l=$R$] (6,3);
\draw (6,3) to[L,l=$L$] (6,0);
\draw (6,0) -- (0,0);
\draw (5,1.5) node[right] {$i(t)$};
\end{tikzpicture}
\end{center}

\textbf{Elemente:}
\begin{itemize}
    \item $U_{\mathrm{DC}}$: Gleichspannungsquelle
    \item $S$: idealer Schalter (Ein/Aus)
    \item $R$: ohmsche Last
    \item $L$: Induktivität
    \item $i(t)$: Laststrom (Zustandsgrösse)
\end{itemize}

\subsubsection{Zwei Topologien pro Periode}

Bei periodischem Schalten entstehen zwei lineare Teilsysteme:

\medskip
\textbf{Topologie 1: Einschaltphase ($0 < t < T_{\mathrm{ein}}$)}

\begin{center}
\begin{tikzpicture}[scale=1]
\draw (0,0) to[battery,l=$U_{\mathrm{DC}}$] (0,3);
\draw (0,3) -- (4,3);
\draw (4,3) to[R,l=$R$] (6,3);
\draw (6,3) to[L,l=$L$] (6,0);
\draw (6,0) -- (0,0);
\end{tikzpicture}
\end{center}

Schalter geschlossen, Quelle aktiv.

\medskip
\textbf{Topologie 2: Ausschaltphase ($T_{\mathrm{ein}} < t < T_s$)}

\begin{center}
\begin{tikzpicture}[scale=1]
\draw (4,3) to[R,l=$R$] (6,3);
\draw (6,3) to[L,l=$L$] (6,0);
\draw (6,0) -- (4,0);
\draw (4,0) -- (4,3);
\draw (3,1.5) node[left] {Freilaufpfad};
\end{tikzpicture}
\end{center}

Quelle getrennt, Strom fliesst im geschlossenen $R$-$L$-Kreis weiter.

\subsection{Differentialgleichungen pro Intervall}

\subsubsection{Einschaltphase}

Kirchhoff:
\[
\cbl{U_{\mathrm{DC}} = R\,i(t) + L\,\frac{di(t)}{dt}}
\]

Standardform:
\[
\cbl{\frac{di}{dt} + \frac{R}{L} i = \frac{U_{\mathrm{DC}}}{L}}
\]

Zeitkonstante:
\[
\cbl{\tau = \frac{L}{R}}
\]

\subsubsection{Ausschaltphase}

Quelle getrennt:
\[
\cbl{0 = R\,i(t) + L\,\frac{di(t)}{dt}}
\]

Standardform:
\[
\cbl{\frac{di}{dt} + \frac{R}{L} i = 0}
\]

\subsection{Lösung der Einschaltphase}

Allgemeine Lösung:
\[
\cbl{i_{\mathrm{ein}}(t) = i_\infty + \big(i(0) - i_\infty\big)\,e^{-t/\tau}},
\]
mit Endwert:
\[
\cbl{i_\infty = \frac{U_{\mathrm{DC}}}{R}}
\]

Interpretation:
\begin{itemize}
    \item Strom steigt exponentiell gegen $U_{\mathrm{DC}}/R$.
    \item Steigung zu Beginn maximal.
    \item Dynamik vollständig durch $\tau = L/R$ bestimmt.
\end{itemize}

\subsection{Lösung der Ausschaltphase}

Homogene Lösung:
\[
\cbl{i_{\mathrm{aus}}(t) = i(T_{\mathrm{ein}})\, e^{-(t-T_{\mathrm{ein}})/\tau}}
\]

Interpretation:
\begin{itemize}
    \item Strom fällt exponentiell ab.
    \item Energie wird in $R$ dissipiert.
    \item Kein neuer Energieeintrag.
\end{itemize}

\subsection{Kontinuitätsbedingungen am Schaltzeitpunkt}

\crd{\textrightarrow\ Induktorstrom ist stetig. Sprünge sind im idealen Modell nicht erlaubt.}


Zentrale Regel:
\[
\cbl{i_L(t^-)=i_L(t^+)}
\]

Konkret:
\begin{itemize}
    \item Anfang Einschaltphase: $i_{\mathrm{ein}}(0)=i_{\min}$
    \item Übergang: $i_{\mathrm{aus}}(T_{\mathrm{ein}})=i_{\mathrm{ein}}(T_{\mathrm{ein}})=i_{\max}$
    \item Periodenende: $i_{\mathrm{aus}}(T_s)=i_{\min}$
\end{itemize}

Diese drei Bedingungen schliessen das System.

\subsection{Stationärer periodischer Betrieb}

Definition:
\[
\cbl{i(t)=i(t+T_s)}
\]

Konkret:
\[
i_{\mathrm{aus}}(T_s) = i_{\mathrm{ein}}(0).
\]

Damit ergeben sich zwei Gleichungen:
\[
\cbl{i_{\max} = i_\infty + (i_{\min}-i_\infty)e^{-T_{\mathrm{ein}}/\tau}}
\]
\[
\cbl{i_{\min} = i_{\max} e^{-T_{\mathrm{aus}}/\tau}}
\]
mit
\[
T_s = T_{\mathrm{ein}} + T_{\mathrm{aus}}.
\]

Diese Gleichungen liefern $i_{\min}$ und $i_{\max}$ explizit.

\subsection{Typische Ergebnisgrössen}

\subsubsection{Stromrippel}
\[
\cbl{\Delta i = i_{\max} - i_{\min}}
\]

\subsubsection{Mittelwert des Stroms}
\[
\overline{i} = \frac{1}{T_s} \left(
\int_0^{T_{\mathrm{ein}}} i_{\mathrm{ein}}(t)\,dt
+
\int_{T_{\mathrm{ein}}}^{T_s} i_{\mathrm{aus}}(t)\,dt
\right).
\]

\subsubsection{Effektivwert des Stroms}
\[
I_{\mathrm{RMS}} = \sqrt{
\frac{1}{T_s}
\left(
\int_0^{T_{\mathrm{ein}}} i_{\mathrm{ein}}^2(t)\,dt
+
\int_{T_{\mathrm{ein}}}^{T_s} i_{\mathrm{aus}}^2(t)\,dt
\right)
}.
\]

\subsubsection{Lastspannungen}

\begin{itemize}
    \item Ohmsch:
    \[
    \cbl{u_R(t) = R\,i(t)}
    \]
    \item Induktiv:
    \[
    \cbl{u_L(t) = L\,\frac{di(t)}{dt}}
    \]
\end{itemize}

Während:
\begin{itemize}
    \item Einschalten: $u_L = U_{\mathrm{DC}} - R i(t)$
    \item Ausschalten: $u_L = -R i(t)$
\end{itemize}

\subsection{Zeitverläufe (qualitativ)}

\begin{center}
\begin{tikzpicture}[scale=1]
\draw[->] (0,0) -- (8,0) node[right] {$t$};
\draw[->] (0,0) -- (0,3) node[above] {$i(t)$};

\draw (0,1) .. controls (1,2.2) .. (3,2.5);
\draw (3,2.5) .. controls (4,2) .. (6,1.2);
\draw (6,1.2) .. controls (7,1.8) .. (8,2.2);

\draw[dashed] (3,0) -- (3,3) node[above] {$T_{\mathrm{ein}}$};
\draw[dashed] (6,0) -- (6,3) node[above] {$T_s$};

\draw (0,1) node[left] {$i_{\min}$};
\draw (3,2.5) node[left] {$i_{\max}$};
\end{tikzpicture}
\end{center}

\textbf{Interpretation:}
\begin{itemize}
    \item Exponentieller Anstieg in Einschaltphase
    \item Exponentieller Abfall in Ausschaltphase
    \item Periodisch wiederholbarer Sägezahn mit Krümmung
\end{itemize}

\subsection{Allgemeiner R--L--C-Zweig im Zeitbereich}

Grundgleichungen:
\[
u(t) = R i(t) + L \frac{di}{dt} + u_C(t)
\]
\[
i_C(t) = C \frac{du_C}{dt}
\]

Zustände:
\[
x(t) = \begin{bmatrix} i(t) \\ u_C(t) \end{bmatrix}
\]

Sonderfälle:

R--L:
\[
u(t) = R i + L \frac{di}{dt},
\qquad \tau = \frac{L}{R}
\]

R--C:
\[
\frac{u}{R} + C \frac{du}{dt} = 0,
\qquad \tau = R C
\]

Allgemeine Lösung (1. Ordnung):
\[
x(t) = x_\infty + (x(t_0)-x_\infty)e^{-(t-t_0)/\tau}
\]

\subsection{Einfluss der Parameter}

\subsubsection{Induktivität $L$}
\begin{itemize}
    \item Größeres $L$ $\Rightarrow$ größeres $\tau$ $\Rightarrow$ kleinerer Stromrippel
    \item Trägere Dynamik
\end{itemize}

\subsubsection{Schaltfrequenz $f_s = 1/T_s$}
\begin{itemize}
    \item Höheres $f_s$ $\Rightarrow$ kleineres $T_s$ $\Rightarrow$ kleinerer Rippel
    \item Höhere Schaltverluste
\end{itemize}

\subsubsection{Widerstand $R$}
\begin{itemize}
    \item Größeres $R$ $\Rightarrow$ kleinerer Endstrom $U_{\mathrm{DC}}/R$
    \item Kürzere Zeitkonstante $\tau$
\end{itemize}

\subsection{Typische Fehlerquellen}

\crd{Diese Fehler kosten in der Prüfung direkt Punkte.}


\begin{itemize}
    \item Falsches Vorzeichen im Exponenten
    \item Falsche Zeitverschiebung in der Ausschaltphase ($t-T_{\mathrm{ein}}$ vergessen)
    \item Kontinuität am Schaltpunkt nicht eingehalten
    \item Periodenbedingung nicht geschlossen
    \item $i_\infty$ falsch angesetzt
\end{itemize}

\subsection{Standard-Vorgehen für Prüfungsaufgaben}

\begin{enumerate}
    \item Schaltintervalle identifizieren.
    \item Pro Intervall DGL aufstellen.
    \item Zeitkonstante $\tau=L/R$ bestimmen.
    \item Lösungen ansetzen (exponentiell).
    \item Kontinuitätsbedingungen formulieren.
    \item Periodenbedingung $i(0)=i(T_s)$ anwenden.
    \item Gesuchte Grössen berechnen: $i_{\min}$, $i_{\max}$, $\Delta i$, $\overline{i}$, $I_{\mathrm{RMS}}$.
\end{enumerate}
